% =====================================================================
% Ensaio 1 — Seção 1: Introdução
%
% Escrita pela fórmula de Head (via Bellemare), na ordem que o framework fixa
% (docs/framework-dissertacao.md §3.1):
%   1. Hook — pelo mundo real, NUNCA pela literatura
%   2. Pergunta de pesquisa — como PRIMEIRA FRASE do parágrafo
%   3. Antecedentes — os estudos mais próximos, contados como história
%   4. Value added — (i) identificação, (ii) validade externa, (iii) mecanismo
%   5. Roadmap
%
% ⚠️ O pecado capital, segundo Bellemare, é o bait-and-switch: prometer na
% introdução o que o corpo não entrega. Como o alvo da qualificação NÃO inclui
% resultado estimado (docs/ars/04-roteiro-qualificacao.md), esta introdução
% promete um DESENHO e o §\ref{sub:promessa} declara isso ao leitor. Nenhuma
% frase abaixo anuncia achado.
%
% ⚠️ Esta seção deve ser REESCRITA a cada iteração do documento — Bellemare
% estima que ~75% do destino de um artigo se decide na introdução.
%
% ✅ REFERÊNCIAS. Diferentemente das §2 e §3 quando foram escritas, esta seção
% usa \cite: paper/referencias.bib foi conferido contra api.crossref.org em
% 2026-08-25 (log em docs/referencias-verificadas.md).
% =====================================================================

\section{Introdução}
\label{sec:introducao}

Em 2010, José Maria Filho --- Zé Maria do Tomé --- foi assassinado na Chapada do
Apodi, no interior do Ceará.\footnote{\textbf{[a conferir]} Circunstâncias do
crime (data exata e número de disparos) \textbf{não} conferidas contra fonte
primária; o texto registra apenas o ano, que consta da tramitação legislativa.
Conferir antes da versão final.} Ele vinha denunciando havia anos a pulverização aérea de
agrotóxicos sobre os perímetros irrigados da região, onde aviões sobrevoam
plantações de banana e melão que fazem fronteira com casas, escolas e açudes de
abastecimento. A denúncia não era sobre o veneno aplicado na lavoura; era sobre o
veneno que \emph{não} fica na lavoura. Nove anos depois, o estado do Ceará
proibiu a prática em todo o território --- a primeira proibição estadual do
gênero no Brasil --- e deu à lei o nome do homem assassinado.

A pulverização aérea é um caso raro de externalidade em que o mecanismo físico é
visível a olho nu. Parte da calda liberada pela aeronave não atinge o alvo: é
carregada pelo vento e se deposita fora da área tratada, num fenômeno que a
literatura técnica chama de deriva. Quem contrata a aplicação paga pelo serviço e
pelo produto; quem mora a favor do vento paga o resto. Essa é a pergunta
distributiva que a lei cearense tenta responder por decreto --- e é uma pergunta
sobre a qual, no Brasil, praticamente não existe evidência causal.

\textbf{Qual é o efeito causal da proibição da pulverização aérea de agrotóxicos
sobre desfechos de saúde ao nascer, e como esse efeito varia com a intensidade
com que cada município era exposto à prática antes da proibição?} A segunda
metade da pergunta não é um refinamento: é o objeto. Um número único --- o efeito
médio da proibição --- diria se a política funcionou, mas não diria se proibir
era o instrumento certo. Essa segunda pergunta, que é a do Ensaio 2 desta
dissertação, se decide pela \emph{inclinação} da curva de dano marginal, e não
pelo seu nível \citep{weitzman1974}. Por isso o parâmetro-alvo aqui é a
\textbf{curva dose-resposta}, e não o efeito num ponto.

\subsection{Antecedentes}
\label{sub:antecedentes}

A evidência econômica mais próxima deste caso se organiza em torno de duas
perguntas que raramente se encontram: o que uma \emph{substância} faz à saúde, e
o que um \emph{método de aplicação} faz à saúde.

Do lado da substância, \citet{dias2023down} são a referência brasileira. Eles
exploram a introdução da soja transgênica no Brasil --- e o ganho potencial de
produtividade que ela trouxe a cada município, que é a variação de aptidão ---
combinada à \emph{direção do fluxo d'água dentro da bacia}, e documentam
deterioração dos desfechos ao nascer nas populações situadas a jusante. Na
especificação preferida, o aumento médio do uso de glifosato entre 2000 e 2010
eleva a mortalidade infantil em 5\% da média amostral.

\citet{reynier2025glyphosate} replicam a lógica nos Estados Unidos, sobre a
adoção de sementes geneticamente modificadas e com o mesmo instrumento de
aptidão, e encontram redução significativa do peso ao nascer e da duração da
gestação --- com efeito \textbf{doze vezes maior} no decil inferior de peso
esperado do que no superior. O dano não é apenas médio; é desigual.
\citet{larsen2017agricultural} acrescentam a forma funcional: no Vale de San
Joaquin, a exposição residencial eleva desfechos adversos em 5 a 9\%, mas
\emph{apenas} entre os expostos ao topo da distribuição --- o percentil 95 para
cima. É resultado que recomenda tratar a dose como contínua e desconfiar de
linearidade. E \citet{frank2024ecosystem} fecham o argumento pelo avesso, com um
choque que empurra o uso de inseticida para \emph{cima}: a perda de morcegos
predadores nos Estados Unidos leva o agricultor a compensar com 31,1\% mais
inseticida, e a mortalidade infantil sobe 7,9\% nos condados atingidos.

Do lado do método, o antecedente direto é \citet{camacho2017health}, que avaliam
as consequências sanitárias da aspersão aérea de herbicida sobre cultivos
ilícitos na Colômbia.\footnote{\textbf{[a conferir]} O registro Crossref deste
artigo não traz \emph{abstract} depositado, de modo que a descrição acima se
apoia no título. Os desfechos específicos e a magnitude serão conferidos contra
o texto integral antes da versão final.} É o estudo mais próximo do objeto
aqui, com uma diferença que importa: na Colômbia a pulverização aérea foi \emph{introduzida} como política de
Estado, ao passo que no Ceará ela foi \emph{retirada}. Os sinais esperados se
invertem, e com eles a estrutura de quem escolhe ser exposto.

O que nenhum desses trabalhos oferece é a avaliação de uma \textbf{proibição de
método} --- o instrumento que a lei cearense de fato usa. Do lado normativo, a
comparação entre proibir e taxar é antiga na economia ambiental
\citep{weitzman1974,lichtenberg1988,helfand1995}, e há simulação recente do
banimento de uma molécula na Europa \citep{bocker2020glyphosate}, mas simulação
ex-ante não é avaliação ex-post. Do lado do método econométrico, a peça que
faltava chegou agora: \citet{callaway2024continuous} mostram que
diferenças-em-diferenças com tratamento contínuo não recupera a curva
dose-resposta sob as hipóteses habituais, e propõem os estimandos que a
recuperam --- distinguindo, em particular, o efeito no nível de dose recebido da
curva causal propriamente dita.

\subsection{Contribuição}
\label{sub:contribuicao}

Este ensaio se situa no cruzamento desses dois lados, e a contribuição tem três
camadas.

\paragraph{Identificação.} A lei cearense entrou em vigor de uma só vez em todo o
estado, de modo que não há variação de \emph{timing} a explorar e a maquinaria de
diferenças-em-diferenças escalonado não se aplica. A identificação vem da
\textbf{intensidade pré-ban de exposição} entre municípios, tratada como dose
contínua, com o estimador de \citet{callaway2024continuous}. O desenho separa
explicitamente o parâmetro que sai sob tendências paralelas usuais daquele que
exige a hipótese forte, e declara --- antes de estimar --- que o teste de
tendências pré-tratamento \emph{não} valida o segundo. Declara também seu próprio
poder: com poucos municípios de dose alta e desfechos ruidosos, o efeito mínimo
detectável é calculado \emph{ex ante}, e coeficiente abaixo dele é reportado como
limite superior, não como achado.

\paragraph{Validade externa.} O objeto é uma \textbf{proibição de método},
raramente avaliada. A distinção não é semântica: proibir o avião não proíbe a
molécula, o que significa que a política atua sobre a \emph{via de exposição}
mantendo a disponibilidade do produto --- e que a aplicação pode migrar para o
solo em vez de desaparecer. A relevância para política é imediata: o Supremo
Tribunal Federal julgou a lei constitucional em 2023 (ADI 6137), de modo que a
competência estadual para legislar nesta matéria está assentada e o instrumento é
replicável por outras unidades da federação.\footnote{\textbf{[a conferir]} A
afirmação de que proposições semelhantes já tramitam em outros estados consta
do material de
projeto, mas \textbf{não} foi conferida contra as casas legislativas
correspondentes, e por isso não é feita aqui.}

\paragraph{Mecanismo.} O desenho não trata o efeito como caixa-preta. Separa o
canal hídrico do canal atmosférico, e acrescenta um terceiro de \emph{sinal
oposto}: se a aplicação migra do avião para o trator e o pulverizador costal, a
proibição afasta o agrotóxico da população e o aproxima do aplicador, de modo que
a intoxicação aguda ocupacional deve \emph{subir}. O mesmo registro hospitalar
fornece a falsificação: a intoxicação autoprovocada responde à disponibilidade da
molécula, que a lei não altera, e portanto \textbf{não} deve se mover com a dose.
Duas séries, mesma população, mesmo sistema de registro, predições opostas.

\paragraph{Uma correção de escopo que a literatura de referência impõe.} Os
trabalhos citados acima são construídos, quase sem exceção, sobre \emph{glifosato
em culturas transgênicas tolerantes a herbicida}. A química documentada na
Chapada do Apodi é outra: a tabela de análises de água reproduzida na
justificativa do projeto de lei registra procimidona e carbaril em 23 de 23
pontos de coleta, carbofurano em 18 e fenitrotiona em 16 --- contra glifosato em
4. É fungicida e inseticida sobre fruticultura irrigada, não herbicida sobre
grão. A consequência é declarada desde já: os desenhos de referência transferem
na \emph{forma} --- instrumentar exposição por aptidão agroclimática ---, não no
\emph{mecanismo}, e a §\ref{sub:quimico} detalha o que isso obriga a mudar.

\subsection{O que este documento entrega, e o que não entrega}
\label{sub:promessa}

\begin{quote}
Este é um documento de \textbf{qualificação}, e seu alvo declarado é um projeto
defensável: identificação fechada, ameaças mapeadas, decisões de especificação
registradas com antecedência. \textbf{Não há resultado estimado aqui}, e a
ausência é deliberada. As seções de resultado estão fora do alvo desta etapa ---
não atrasadas em relação a ele.
\end{quote}

A escolha tem uma razão que vai além do cronograma. O conjunto de desfechos que
os registros administrativos brasileiros tornam disponíveis é amplo, e o poder
estatístico deste desenho é curto; essa combinação é precisamente a que produz
achado espúrio com aparência de rigor. Fixar desfecho primário, exposição
primária e critério de falsificação \emph{antes} do primeiro contato com dado
real é, aqui, condição de credibilidade --- e é verificável, porque o registro é
versionado.

\subsection{Organização do trabalho}
\label{sub:roadmap}

O restante deste ensaio organiza-se como segue. A §\ref{sec:background} apresenta
a lei, sua cronologia --- que distingue três marcos de antecipação, e não um ---,
a região a que ela responde e o mecanismo físico da deriva, além da correção
química discutida acima. A §\ref{sec:teoria} desenvolve o referencial de
externalidade e escolha de instrumento, e termina no conjunto de predições
testáveis que amarra a teoria à estimação. A §\ref{sec:identificacao} ---
o núcleo deste documento --- expõe o parâmetro-alvo, as hipóteses que ele exige,
a construção da variável de dose e do grupo de comparação, o poder do desenho, os
procedimentos de inferência e a tabela de ameaças, cada uma com o que faz ao
estimando e o que o desenho lhe opõe. As seções de dados, resultados e conclusão
completam o ensaio na versão final da dissertação.
